\documentclass[acmsmall,screen]{acmart}
\pagestyle{plain}

\usepackage{tikz}
\usetikzlibrary{matrix,fit,quantikz}
\usepackage{pgfplots}
\usepgfplotslibrary{groupplots,colorbrewer}
\newcommand{\expand}{\mathop{\mathsf{expand}}}

%%
%% end of the preamble, start of the body of the document source.
\begin{document}
\begin{figure}
\pgfplotsset{
  small,
  compat=1.5,
  every tick label/.append style={font=\tiny},
  label style={font=\footnotesize},
  title style={font=\footnotesize},
  ymode=log,
  table/col sep=comma,
  colormap/Paired-10,
  % styles
  fggi/.style={
    color=Set1-A,
    mark=*,
    solid,
    mark size=1.8pt,
    mark options={draw=Paired-F},
  },
  fggp/.style={
    color=Set1-B,
    mark=square*,
    solid,
    mark size=1.8pt,
    mark options={draw=Paired-B,rotate=45},
  },
  qiskit/.style={
    color=Set1-C,
    mark=Mercedes star,
    solid,
    mark size=1.8pt,
  },
  cirq/.style={
    color=Set1-D,
    mark=x,
    solid,
    mark size=1.8pt,
  },
  qtorch/.style={
    color=Set1-E,
    mark=star,
    solid,
    mark size=1.8pt,
  },
  quimb/.style={
    color=Set1-H,
    mark=pentagon,
    solid,
    mark size=1.8pt,
  },
  gates/.style={
    color=gray,
    mark=none,
    dotted,
    semithick,
  },
}
\begin{center}
\begin{tikzpicture}
  \begin{groupplot}[group style={group size=2 by 3, group name=b1},
    group/horizontal sep=4em,
    group/vertical sep=5em]

    \nextgroupplot[
    xlabel={Number of qubits},
    ylabel={Running time (seconds)},
    axis y line*=left,
    title={Hadamard to the last qubit},
    legend style={
      at={(0.5,-0.1)},
      anchor=north,legend columns=-1,
      font=\footnotesize,
      draw=none,
      legend columns=2,
      legend cell align={left},
    },
    legend to name=benchmark1,
    mark repeat=8,
    execute at end plot visualization/.append={%
      \xdef\savedxminhad{\pgfkeysvalueof{/pgfplots/xmin}}%
      \xdef\savedxmaxhad{\pgfkeysvalueof{/pgfplots/xmax}}%
      \xdef\savedyminhad{\pgfkeysvalueof{/pgfplots/ymin}}%
      \xdef\savedymaxhad{\pgfkeysvalueof{/pgfplots/ymax}}%
    },
    ]
    \addplot[fggi] table[x=qubit-number, y=had-last-qubit-allein_opt:iso-fgg-total] {had-last-qubit.csv};
    \addplot[fggp, mark phase=3] table[x=qubit-number, y=had-last-qubit-allein_opt:iso-perpl-fgg-total] {had-last-qubit.csv};
    \addplot[qiskit] table[x=qubit-number, y=had-last-qubit-avg:qiskit-simulation] {had-last-qubit.csv};
    \addplot[cirq] table[x=qubit-number, y=had-last-qubit-avg:qsim-simulation] {had-last-qubit.csv};
    \addplot[quimb] table[x=qubit-number, y=had-last-qubit-avg:quimb-simulation] {had-last-qubit.csv};

    \addlegendentry{FGG simulation with $\expand$ and $\bigl[\begin{smallmatrix}1&0\\0&1\end{smallmatrix}\bigr]$ opts};
    \addlegendentry{FGG simulation with $\expand$ and $\bigl[\begin{smallmatrix}1&0\\0&1\end{smallmatrix}\bigr]$ opts, through PERPL};
    \addlegendentry{Qiskit};
    \addlegendentry{qsim};
    \addlegendentry{quimb};
    \addlegendimage{qtorch};
    \addlegendentry{qTorch};
    \addlegendimage{gates};
    \addlegendentry{Number of gates};

    \nextgroupplot[
    xlabel={Number of qubits},
    axis y line*=left,
    title={Deutsch-Jozsa with the even oracle},
    mark repeat=4,
    execute at end plot visualization/.append={%
      \xdef\savedxmindje{\pgfkeysvalueof{/pgfplots/xmin}}%
      \xdef\savedxmaxdje{\pgfkeysvalueof{/pgfplots/xmax}}%
      \xdef\savedymindje{\pgfkeysvalueof{/pgfplots/ymin}}%
      \xdef\savedymaxdje{\pgfkeysvalueof{/pgfplots/ymax}}%
    },
    ]
    \addplot[fggi] table[x=qubit-number, y=deutsch-jozsa-is-even-simplified-allein_opt:iso-fgg-total] {deutsch-jozsa-is-even-simplified.csv};
    \addplot[fggp, mark phase=3] table[x=qubit-number, y=deutsch-jozsa-is-even-simplified-allein_opt:iso-perpl-fgg-total] {deutsch-jozsa-is-even-simplified.csv};
    \addplot[qiskit] table[x=qubit-number, y=deutsch-jozsa-is-even-simplified-avg:qiskit-simulation] {deutsch-jozsa-is-even-simplified.csv};
    \addplot[cirq] table[x=qubit-number, y=deutsch-jozsa-is-even-simplified-avg:qsim-simulation] {deutsch-jozsa-is-even-simplified.csv};
    \addplot[quimb] table[x=qubit-number, y=deutsch-jozsa-is-even-simplified-avg:quimb-simulation] {deutsch-jozsa-is-even-simplified.csv};

    \nextgroupplot[
    xlabel={Number of qubits},
    ylabel={Running time (seconds)},
    axis y line*=left,
    title={Simon's},
    mark repeat=1,
    execute at end plot visualization/.append={%
      \xdef\savedxminsimon{\pgfkeysvalueof{/pgfplots/xmin}}%
      \xdef\savedxmaxsimon{\pgfkeysvalueof{/pgfplots/xmax}}%
      \xdef\savedyminsimon{\pgfkeysvalueof{/pgfplots/ymin}}%
      \xdef\savedymaxsimon{\pgfkeysvalueof{/pgfplots/ymax}}%
    },
    ]
    \addplot[fggi] table[x=qubit-number, y=simon-allein_opt:iso-fgg-total] {simon.csv};
    \addplot[fggp] table[x=qubit-number, y=simon-allein_opt:iso-perpl-fgg-total] {simon.csv};
    \addplot[qiskit] table[x=qubit-number, y=simon-avg:qiskit-simulation] {simon.csv};
    \addplot[cirq] table[x=qubit-number, y=simon-avg:qsim-simulation] {simon.csv};
    \addplot[quimb] table[x=qubit-number, y=simon-avg:quimb-simulation] {simon.csv};

    \nextgroupplot[
    xlabel={Number of qubits},
    axis y line*=left,
    title={Decomposed $C^{n}(X)$},
    mark repeat=2,
    execute at end plot visualization/.append={%
      \xdef\savedxminmcx{\pgfkeysvalueof{/pgfplots/xmin}}%
      \xdef\savedxmaxmcx{\pgfkeysvalueof{/pgfplots/xmax}}%
      \xdef\savedyminmcx{\pgfkeysvalueof{/pgfplots/ymin}}%
      \xdef\savedymaxmcx{\pgfkeysvalueof{/pgfplots/ymax}}%
    },
    ]
    \addplot[fggi] table[x=qubit-number, y=mcx-allein_opt:iso-fgg-total] {mcx.csv};
    \addplot[fggp] table[x=qubit-number, y=mcx-allein_opt:iso-perpl-fgg-total] {mcx.csv};
    \addplot[qiskit] table[x=qubit-number, y=mcx-avg:qiskit-simulation] {mcx.csv};
    \addplot[cirq] table[x=qubit-number, y=mcx-avg:qsim-simulation] {mcx.csv};
    \addplot[quimb] table[x=qubit-number, y=mcx-avg:quimb-simulation] {mcx.csv};

    \nextgroupplot[
    xlabel={Number of qubits},
    ylabel={Running time (seconds)},
    axis y line*=left,
    title={QFT},
    mark repeat=2,
    execute at end plot visualization/.append={%
      \xdef\savedxminqft{\pgfkeysvalueof{/pgfplots/xmin}}%
      \xdef\savedxmaxqft{\pgfkeysvalueof{/pgfplots/xmax}}%
      \xdef\savedyminqft{\pgfkeysvalueof{/pgfplots/ymin}}%
      \xdef\savedymaxqft{\pgfkeysvalueof{/pgfplots/ymax}}%
    },
    ]
    \addplot[fggi] table[x=qubit-number, y=qft-allein_opt:iso-fgg-total] {qft.csv};
    \addplot[fggp] table[x=qubit-number, y=qft-allein_opt:iso-perpl-fgg-total] {qft.csv};
    \addplot[qiskit] table[x=qubit-number, y=qft-avg:qiskit-simulation] {qft.csv};
    \addplot[cirq] table[x=qubit-number, y=qft-avg:qsim-simulation] {qft.csv};
    \addplot[quimb] table[x=qubit-number, y=qft-avg:quimb-simulation] {qft.csv};

    \nextgroupplot[
    xlabel={Number of qubits},
    axis y line*=left,
    title={Grover},
    mark repeat=2,
    execute at end plot visualization/.append={%
      \xdef\savedxmingrover{\pgfkeysvalueof{/pgfplots/xmin}}%
      \xdef\savedxmaxgrover{\pgfkeysvalueof{/pgfplots/xmax}}%
      \xdef\savedymingrover{\pgfkeysvalueof{/pgfplots/ymin}}%
      \xdef\savedymaxgrover{\pgfkeysvalueof{/pgfplots/ymax}}%
    },
    ]
    \addplot[fggi] table[x=qubit-number, y=grover-allein_opt:iso-fgg-total] {grover.csv};
    \addplot[fggp] table[x=qubit-number, y=grover-allein_opt:iso-perpl-fgg-total] {grover.csv};
    \addplot[qiskit] table[x=qubit-number, y=grover-avg:qiskit-simulation] {grover.csv};
    \addplot[cirq] table[x=qubit-number, y=grover-avg:qsim-simulation] {grover.csv};
    \addplot[quimb] table[x=qubit-number, y=grover-avg:quimb-simulation] {grover.csv};
  \end{groupplot}

  \begin{axis}[
    at={(b1 c1r1.south west)},
    anchor=south west,
    ymode=log,
    xmin=\savedxminhad,
    xmax=\savedxmaxhad,
    ymin=exp(\savedyminhad),
    ymax=exp(\savedymaxhad),
    mark repeat=4,
    hide axis,
    ]
    \addplot[qtorch] table[x=qubit-number, y=had-last-qubit-avg:qtorch-simulation-full] {had-last-qubit.csv};
  \end{axis}

  \begin{axis}[
    at={(b1 c2r1.south west)},
    anchor=south west,
    ymode=log,
    xmin=\savedxmindje,
    xmax=\savedxmaxdje,
    ymin=exp(\savedymindje),
    ymax=exp(\savedymaxdje),
    mark repeat=2,
    hide axis,
    ]
    \addplot[qtorch] table[x=qubit-number, y=deutsch-jozsa-is-even-simplified-avg:qtorch-simulation-full] {deutsch-jozsa-is-even-simplified.csv};
  \end{axis}

  \begin{axis}[
    at={(b1 c1r2.south west)},
    anchor=south west,
    ymode=log,
    xmin=\savedxminsimon,
    xmax=\savedxmaxsimon,
    ymin=exp(\savedyminsimon),
    ymax=exp(\savedymaxsimon),
    hide axis,
    ]
    \addplot[qtorch] table[x=qubit-number, y=simon-avg:qtorch-simulation-full] {simon.csv};
  \end{axis}

  \begin{axis}[
    at={(b1 c2r2.south west)},
    anchor=south west,
    ymode=log,
    xmin=\savedxminmcx,
    xmax=\savedxmaxmcx,
    ymin=exp(\savedyminmcx),
    ymax=exp(\savedymaxmcx),
    hide axis,
    ]
    \addplot[qtorch] table[x=qubit-number, y=mcx-avg:qtorch-simulation-full] {mcx.csv};
  \end{axis}

  \begin{axis}[
    at={(b1 c1r3.south west)},
    anchor=south west,
    ymode=log,
    xmin=\savedxminqft,
    xmax=\savedxmaxqft,
    ymin=exp(\savedyminqft),
    ymax=exp(\savedymaxqft),
    hide axis,
    ]
    \addplot[qtorch] table[x=qubit-number, y=qft-avg:qtorch-simulation-full] {qft.csv};
  \end{axis}

  \begin{axis}[
    at={(b1 c2r3.south west)},
    anchor=south west,
    ymode=log,
    xmin=\savedxmingrover,
    xmax=\savedxmaxgrover,
    ymin=exp(\savedymingrover),
    ymax=exp(\savedymaxgrover),
    hide axis,
    ]
    \addplot[qtorch] table[x=qubit-number, y=grover-avg:qtorch-simulation-full] {grover.csv};
  \end{axis}

  \begin{axis}[
    at={(b1 c1r1.south west)},
    anchor=south west,
    axis x line=none,
    axis y line*=right,
    ymode=log,
    ]
    \addplot[gates] table[x=qubit-number, y=gate-number] {had-last-qubit.csv};
  \end{axis}

  \begin{axis}[
    at={(b1 c2r1.south west)},
    anchor=south west,
    axis x line=none,
    axis y line*=right,
    ymode=log,
    ylabel={Number of gates},
    ]
    \addplot[gates] table[x=qubit-number, y=gate-number] {deutsch-jozsa-is-even-simplified.csv};
  \end{axis}

  \begin{axis}[
    at={(b1 c1r2.south west)},
    anchor=south west,
    axis x line=none,
    axis y line*=right,
    ymode=log,
    ]
    \addplot[gates] table[x=qubit-number, y=gate-number] {simon.csv};
  \end{axis}

  \begin{axis}[
    at={(b1 c2r2.south west)},
    anchor=south west,
    axis x line=none,
    axis y line*=right,
    ymode=log,
    ylabel={Number of gates},
    ]
    \addplot[gates] table[x=qubit-number, y=gate-number] {mcx.csv};
  \end{axis}

  \begin{axis}[
    at={(b1 c1r3.south west)},
    anchor=south west,
    axis x line=none,
    axis y line*=right,
    ymode=log,
    ]
    \addplot[gates] table[x=qubit-number, y=gate-number] {qft.csv};
  \end{axis}

  \begin{axis}[
    at={(b1 c2r3.south west)},
    anchor=south west,
    axis x line=none,
    axis y line*=right,
    ymode=log,
    ylabel={Number of gates},
    ]
    \addplot[gates] table[x=qubit-number, y=gate-number] {grover.csv};
  \end{axis}

  \path (b1 c1r3.south west) -- (b1 c2r3.south east) coordinate[midway] (lower_middle);
  \node[anchor=north, yshift=-2.5em] at (lower_middle) {\ref{benchmark1}};
\end{tikzpicture}
\end{center}
\caption{Simulation and compilation times of our FGG-based simulators, comparing different sets of optimizations.
  For each plot, the $x$-axis shows the number of qubits, and the $y$-axis shows run time in seconds.}
\label{fig:bench-1}
\end{figure}

\newpage
\pagestyle{plain}
\begin{figure}
\pgfplotsset{
  small,
  compat=1.5,
  every tick label/.append style={font=\tiny},
  label style={font=\footnotesize},
  title style={font=\footnotesize},
  ymode=log,
  table/col sep=comma,
  colormap/Paired-10,
  % styles
  fggi/.style={
    color=Set1-A,
    mark=*,
    solid,
    mark size=1.8pt,
    mark options={draw=Paired-F},
  },
  fggp/.style={
    color=Set1-B,
    mark=square*,
    solid,
    mark size=1.8pt,
    mark options={draw=Paired-B,rotate=45},
  },
  qiskit/.style={
    color=Set1-C,
    mark=Mercedes star,
    solid,
    mark size=1.8pt,
  },
  cirq/.style={
    color=Set1-D,
    mark=x,
    solid,
    mark size=1.8pt,
  },
  qtorch/.style={
    color=Set1-E,
    mark=star,
    solid,
    mark size=1.8pt,
  },
  quimb/.style={
    color=Set1-H,
    mark=pentagon,
    solid,
    mark size=1.8pt,
  },
  gates/.style={
    color=gray,
    mark=none,
    dotted,
    semithick,
  },
}
\begin{center}
\begin{tikzpicture}
  \begin{groupplot}[group style={group size=2 by 3, group name=b1},
    group/horizontal sep=4em,
    group/vertical sep=5em]
    \nextgroupplot[
    xlabel={Number of qubits},
    ylabel={Running time (seconds)},
    axis y line*=left,
    title={Bell state},
    legend style={
      at={(0.5,-0.1)},
      anchor=north,legend columns=-1,
      font=\footnotesize,
      draw=none,
      legend columns=2,
      legend cell align={left},
    },
    legend to name=benchmark1,
    mark repeat=8,
    execute at end plot visualization/.append={%
      \xdef\savedxminbell{\pgfkeysvalueof{/pgfplots/xmin}}%
      \xdef\savedxmaxbell{\pgfkeysvalueof{/pgfplots/xmax}}%
      \xdef\savedyminbell{\pgfkeysvalueof{/pgfplots/ymin}}%
      \xdef\savedymaxbell{\pgfkeysvalueof{/pgfplots/ymax}}%
    },
    ]
    \addplot[fggi] table[x=qubit-number, y=bell-state-allein_opt:iso-fgg-total] {bell-state.csv};
    \addplot[fggp, mark phase=3] table[x=qubit-number, y=bell-state-allein_opt:iso-perpl-fgg-total] {bell-state.csv};
    \addplot[qiskit] table[x=qubit-number, y=bell-state-avg:qiskit-simulation] {bell-state.csv};
    \addplot[cirq] table[x=qubit-number, y=bell-state-avg:qsim-simulation] {bell-state.csv};
    \addplot[quimb] table[x=qubit-number, y=bell-state-avg:quimb-simulation] {bell-state.csv};

    \addlegendentry{FGG simulation with $\expand$ and $\bigl[\begin{smallmatrix}1&0\\0&1\end{smallmatrix}\bigr]$ opts};
    \addlegendentry{FGG simulation with $\expand$ and $\bigl[\begin{smallmatrix}1&0\\0&1\end{smallmatrix}\bigr]$ opts, through PERPL};
    \addlegendentry{Qiskit};
    \addlegendentry{qsim};
    \addlegendentry{quimb};
    \addlegendimage{qtorch};
    \addlegendentry{qTorch};
    \addlegendimage{gates};
    \addlegendentry{Number of gates};

    \nextgroupplot[
    xlabel={Number of qubits},
    axis y line*=left,
    title={Two $N$-qubit Bell in cascade},
    mark repeat=4,
    execute at end plot visualization/.append={%
      \xdef\savedxminctb{\pgfkeysvalueof{/pgfplots/xmin}}%
      \xdef\savedxmaxctb{\pgfkeysvalueof{/pgfplots/xmax}}%
      \xdef\savedyminctb{\pgfkeysvalueof{/pgfplots/ymin}}%
      \xdef\savedymaxctb{\pgfkeysvalueof{/pgfplots/ymax}}%
    },
    ]
    \addplot[fggi] table[x=qubit-number, y=casc-two-n-bell-state-allein_opt:iso-fgg-total] {casc-two-n-bell-state.csv};
    \addplot[fggp] table[x=qubit-number, y=casc-two-n-bell-state-allein_opt:iso-perpl-fgg-total] {casc-two-n-bell-state.csv};
    \addplot[qiskit] table[x=qubit-number, y=casc-two-n-bell-state-avg:qiskit-simulation] {casc-two-n-bell-state.csv};
    \addplot[cirq] table[x=qubit-number, y=casc-two-n-bell-state-avg:qsim-simulation] {casc-two-n-bell-state.csv};
    \addplot[quimb] table[x=qubit-number, y=casc-two-n-bell-state-avg:quimb-simulation] {casc-two-n-bell-state.csv};

    \nextgroupplot[
    xlabel={Number of qubits},
    ylabel={Running time (seconds)},
    axis y line*=left,
    title={Hadamard to the first qubit before Bell},
    mark repeat=4,
    execute at end plot visualization/.append={%
      \xdef\savedxminchfb{\pgfkeysvalueof{/pgfplots/xmin}}%
      \xdef\savedxmaxchfb{\pgfkeysvalueof{/pgfplots/xmax}}%
      \xdef\savedyminchfb{\pgfkeysvalueof{/pgfplots/ymin}}%
      \xdef\savedymaxchfb{\pgfkeysvalueof{/pgfplots/ymax}}%
    },
    ]
    \addplot[fggi] table[x=qubit-number, y=casc-had-first-bell-allein_opt:iso-fgg-total] {casc-had-first-bell.csv};
    \addplot[fggp, mark phase=3] table[x=qubit-number, y=casc-had-first-bell-allein_opt:iso-perpl-fgg-total] {casc-had-first-bell.csv};
    \addplot[qiskit] table[x=qubit-number, y=casc-had-first-bell-avg:qiskit-simulation] {casc-had-first-bell.csv};
    \addplot[cirq] table[x=qubit-number, y=casc-had-first-bell-avg:qsim-simulation] {casc-had-first-bell.csv};
    \addplot[quimb] table[x=qubit-number, y=casc-had-first-bell-avg:quimb-simulation] {casc-had-first-bell.csv};

    \nextgroupplot[
    xlabel={Number of qubits},
    axis y line*=left,
    title={Hadamard to the last qubit before Bell},
    mark repeat=4,
    execute at end plot visualization/.append={%
      \xdef\savedxminchlb{\pgfkeysvalueof{/pgfplots/xmin}}%
      \xdef\savedxmaxchlb{\pgfkeysvalueof{/pgfplots/xmax}}%
      \xdef\savedyminchlb{\pgfkeysvalueof{/pgfplots/ymin}}%
      \xdef\savedymaxchlb{\pgfkeysvalueof{/pgfplots/ymax}}%
    },
    ]
    \addplot[fggi] table[x=qubit-number, y=casc-had-last-bell-allein_opt:iso-fgg-total] {casc-had-last-bell.csv};
    \addplot[fggp] table[x=qubit-number, y=casc-had-last-bell-allein_opt:iso-perpl-fgg-total] {casc-had-last-bell.csv};
    \addplot[qiskit] table[x=qubit-number, y=casc-had-last-bell-avg:qiskit-simulation] {casc-had-last-bell.csv};
    \addplot[cirq] table[x=qubit-number, y=casc-had-last-bell-avg:qsim-simulation] {casc-had-last-bell.csv};
    \addplot[quimb] table[x=qubit-number, y=casc-had-last-bell-avg:quimb-simulation] {casc-had-last-bell.csv};

    \nextgroupplot[
    xlabel={Number of qubits},
    ylabel={Running time (seconds)},
    axis y line*=left,
    title={One Hadamard in parallel with Deutsch-Jozsa},
    mark repeat=4,
    execute at end plot visualization/.append={%
      \xdef\savedxminhdj{\pgfkeysvalueof{/pgfplots/xmin}}%
      \xdef\savedxmaxhdj{\pgfkeysvalueof{/pgfplots/xmax}}%
      \xdef\savedyminhdj{\pgfkeysvalueof{/pgfplots/ymin}}%
      \xdef\savedymaxhdj{\pgfkeysvalueof{/pgfplots/ymax}}%
    },
    ]
    \addplot[fggi] table[x=qubit-number, y=had-last-dj-even-allein_opt:iso-fgg-total] {had-last-dj-even.csv};
    \addplot[fggp] table[x=qubit-number, y=had-last-dj-even-allein_opt:iso-perpl-fgg-total] {had-last-dj-even.csv};
    \addplot[qiskit] table[x=qubit-number, y=had-last-dj-even-avg:qiskit-simulation] {had-last-dj-even.csv};
    \addplot[cirq] table[x=qubit-number, y=had-last-dj-even-avg:qsim-simulation] {had-last-dj-even.csv};
    \addplot[quimb] table[x=qubit-number, y=had-last-dj-even-avg:quimb-simulation] {had-last-dj-even.csv};
  \end{groupplot}

  \begin{axis}[
    at={(b1 c1r1.south west)},
    anchor=south west,
    ymode=log,
    xmin=\savedxminbell,
    xmax=\savedxmaxbell,
    ymin=exp(\savedyminbell),
    ymax=exp(\savedymaxbell),
    hide axis,
    ]
    \addplot[qtorch] table[x=qubit-number, y=bell-state-avg:qtorch-simulation-full] {bell-state.csv};
  \end{axis}


  \begin{axis}[
    at={(b1 c2r1.south west)},
    anchor=south west,
    ymode=log,
    xmin=\savedxminctb,
    xmax=\savedxmaxctb,
    ymin=exp(\savedymaxctb),
    ymax=exp(\savedymaxctb),
    hide axis,
    ]
    \addplot[qtorch] table[x=qubit-number, y=casc-two-n-bell-state-avg:qtorch-simulation-full] {casc-two-n-bell-state.csv};
  \end{axis}

  \begin{axis}[
    at={(b1 c1r2.south west)},
    anchor=south west,
    ymode=log,
    xmin=\savedxminchfb,
    xmax=\savedxmaxchfb,
    ymin=exp(\savedyminchfb),
    ymax=exp(\savedymaxchfb),
    hide axis,
    ]
    \addplot[qtorch] table[x=qubit-number, y=casc-had-first-bell-avg:qtorch-simulation-full] {casc-had-first-bell.csv};
  \end{axis}

  \begin{axis}[
    at={(b1 c2r2.south west)},
    anchor=south west,
    ymode=log,
    xmin=\savedxminchlb,
    xmax=\savedxmaxchlb,
    ymin=exp(\savedyminchlb),
    ymax=exp(\savedymaxchlb),
    hide axis,
    ]
    \addplot[qtorch] table[x=qubit-number, y=casc-had-last-bell-avg:qtorch-simulation-full] {casc-had-last-bell.csv};
  \end{axis}

  \begin{axis}[
    at={(b1 c1r3.south west)},
    anchor=south west,
    ymode=log,
    xmin=\savedxminhdj,
    xmax=\savedxmaxhdj,
    ymin=exp(\savedyminhdj),
    ymax=exp(\savedymaxhdj),
    hide axis,
    ]
    \addplot[qtorch] table[x=qubit-number, y=had-last-dj-even-avg:qtorch-simulation-full] {had-last-dj-even.csv};
  \end{axis}

  \begin{axis}[
    at={(b1 c1r1.south west)},
    anchor=south west,
    axis x line=none,
    axis y line*=right,
    ymode=log,
    ]
    \addplot[gates] table[x=qubit-number, y=gate-number] {bell-state.csv};
  \end{axis}

  \begin{axis}[
    at={(b1 c2r1.south west)},
    anchor=south west,
    axis x line=none,
    axis y line*=right,
    ymode=log,
    ylabel={Number of gates},
    ]
    \addplot[gates] table[x=qubit-number, y=gate-number] {casc-two-n-bell-state.csv};
  \end{axis}

  \begin{axis}[
    at={(b1 c1r2.south west)},
    anchor=south west,
    axis x line=none,
    axis y line*=right,
    ymode=log,
    ]
    \addplot[gates] table[x=qubit-number, y=gate-number] {casc-had-first-bell.csv};
  \end{axis}

  \begin{axis}[
    at={(b1 c2r2.south west)},
    anchor=south west,
    axis x line=none,
    axis y line*=right,
    ymode=log,
    ylabel={Number of gates},
    ]
    \addplot[gates] table[x=qubit-number, y=gate-number] {casc-had-last-bell.csv};
  \end{axis}

  \begin{axis}[
    at={(b1 c1r3.south west)},
    anchor=south west,
    axis x line=none,
    axis y line*=right,
    ymode=log,
    ylabel={Number of gates},
    ]
    \addplot[gates] table[x=qubit-number, y=gate-number] {had-last-dj-even.csv};
  \end{axis}

  \path (b1 c1r3.south west) -- (b1 c2r3.south east) coordinate[midway] (lower_middle);
  \node[anchor=north, yshift=-2.5em] at (lower_middle) {\ref{benchmark1}};
\end{tikzpicture}
\end{center}
\caption{Simulation and compilation times of our FGG-based simulators, comparing different sets of optimizations.
  For each plot, the $x$-axis shows the number of qubits, and the $y$-axis shows run time in seconds (continued).}
\label{fig:bench-2}
\end{figure}

\newpage
\pagestyle{plain}
\begin{figure}
\pgfplotsset{
  small,
  compat=1.5,
  every tick label/.append style={font=\tiny},
  label style={font=\footnotesize},
  title style={font=\footnotesize},
  ymode=log,
  table/col sep=comma,
  colormap/Paired-10,
  % styles
  fggi/.style={
    color=Set1-A,
    mark=*,
    solid,
    mark size=1.8pt,
    mark options={draw=Paired-F},
  },
  fggp/.style={
    color=Set1-B,
    mark=square*,
    solid,
    mark size=1.8pt,
    mark options={draw=Paired-B,rotate=45},
  },
  qiskit/.style={
    color=Set1-C,
    mark=Mercedes star,
    solid,
    mark size=1.8pt,
  },
  cirq/.style={
    color=Set1-D,
    mark=x,
    solid,
    mark size=1.8pt,
  },
  qtorch/.style={
    color=Set1-E,
    mark=star,
    solid,
    mark size=1.8pt,
  },
  quimb/.style={
    color=Set1-H,
    mark=pentagon,
    solid,
    mark size=1.8pt,
  },
  gates/.style={
    color=gray,
    mark=none,
    dotted,
    semithick,
  },
}
\begin{center}
\begin{tikzpicture}
  \begin{groupplot}[group style={group size=2 by 2, group name=b1},
    group/horizontal sep=4em,
    group/vertical sep=5em]
    \nextgroupplot[
    xlabel={Number of qubits},
    ylabel={Running time (seconds)},
    axis y line*=left,
    title={Hadamard in parallel with $N$-qubit Bell},
    legend style={
      at={(0.5,-0.1)},
      anchor=north,legend columns=-1,
      font=\footnotesize,
      draw=none,
      legend columns=2,
      legend cell align={left},
    },
    legend to name=benchmark1,
    mark repeat=4,
    execute at end plot visualization/.append={%
      \xdef\savedxminhlnb{\pgfkeysvalueof{/pgfplots/xmin}}%
      \xdef\savedxmaxhlnb{\pgfkeysvalueof{/pgfplots/xmax}}%
      \xdef\savedyminhlnb{\pgfkeysvalueof{/pgfplots/ymin}}%
      \xdef\savedymaxhlnb{\pgfkeysvalueof{/pgfplots/ymax}}%
    },
    ]
    \addplot[fggi] table[x=qubit-number, y=para-n-had-last-n-bell-allein_opt:iso-fgg-total] {para-n-had-last-n-bell.csv};
    \addplot[fggp, mark phase=3] table[x=qubit-number, y=para-n-had-last-n-bell-allein_opt:iso-perpl-fgg-total] {para-n-had-last-n-bell.csv};
    \addplot[qiskit] table[x=qubit-number, y=para-n-had-last-n-bell-avg:qiskit-simulation] {para-n-had-last-n-bell.csv};
    \addplot[cirq] table[x=qubit-number, y=para-n-had-last-n-bell-avg:qsim-simulation] {para-n-had-last-n-bell.csv};
    \addplot[quimb] table[x=qubit-number, y=para-n-had-last-n-bell-avg:quimb-simulation] {para-n-had-last-n-bell.csv};

    \addlegendentry{FGG simulation with $\expand$ and $\bigl[\begin{smallmatrix}1&0\\0&1\end{smallmatrix}\bigr]$ opts};
    \addlegendentry{FGG simulation with $\expand$ and $\bigl[\begin{smallmatrix}1&0\\0&1\end{smallmatrix}\bigr]$ opts, through PERPL};
    \addlegendentry{Qiskit};
    \addlegendentry{qsim};
    \addlegendentry{quimb};
    \addlegendimage{qtorch};
    \addlegendentry{qTorch};
    \addlegendimage{gates};
    \addlegendentry{Number of gates};

    \nextgroupplot[
    xlabel={Number of qubits},
    axis y line*=left,
    title={Hadamard in parallel with 10-qubit Bell},
    mark repeat=8,
    execute at end plot visualization/.append={%
      \xdef\savedxminphadtb{\pgfkeysvalueof{/pgfplots/xmin}}%
      \xdef\savedxmaxphadtb{\pgfkeysvalueof{/pgfplots/xmax}}%
      \xdef\savedyminphadtb{\pgfkeysvalueof{/pgfplots/ymin}}%
      \xdef\savedymaxphadtb{\pgfkeysvalueof{/pgfplots/ymax}}%
    },
    ]
    \addplot[fggi] table[x=qubit-number, y=para-n-had-last-ten-qubit-bell-allein_opt:iso-fgg-total] {para-n-had-last-ten-qubit-bell.csv};
    \addplot[fggp] table[x=qubit-number, y=para-n-had-last-ten-qubit-bell-allein_opt:iso-perpl-fgg-total] {para-n-had-last-ten-qubit-bell.csv};
    \addplot[qiskit] table[x=qubit-number, y=para-n-had-last-ten-qubit-bell-avg:qiskit-simulation] {para-n-had-last-ten-qubit-bell.csv};
    \addplot[cirq] table[x=qubit-number, y=para-n-had-last-ten-qubit-bell-avg:qsim-simulation] {para-n-had-last-ten-qubit-bell.csv};
    \addplot[quimb] table[x=qubit-number, y=para-n-had-last-ten-qubit-bell-avg:quimb-simulation] {para-n-had-last-ten-qubit-bell.csv};

    \nextgroupplot[
    xlabel={Number of qubits},
    ylabel={Running time (seconds)},
    axis y line*=left,
    title={$N$ 10-qubit Bell in parallel},
    execute at end plot visualization/.append={%
      \xdef\savedxminntbell{\pgfkeysvalueof{/pgfplots/xmin}}%
      \xdef\savedxmaxntbell{\pgfkeysvalueof{/pgfplots/xmax}}%
      \xdef\savedyminntbell{\pgfkeysvalueof{/pgfplots/ymin}}%
      \xdef\savedymaxntbell{\pgfkeysvalueof{/pgfplots/ymax}}%
    },
    ]
    \addplot[fggi] table[x=qubit-number, y=para-n-ten-qubit-bell-allein_opt:iso-fgg-total] {para-n-ten-qubit-bell.csv};
    \addplot[fggp, mark phase=3] table[x=qubit-number, y=para-n-ten-qubit-bell-allein_opt:iso-perpl-fgg-total] {para-n-ten-qubit-bell.csv};
    \addplot[qiskit] table[x=qubit-number, y=para-n-ten-qubit-bell-avg:qiskit-simulation] {para-n-ten-qubit-bell.csv};
    \addplot[cirq] table[x=qubit-number, y=para-n-ten-qubit-bell-avg:qsim-simulation] {para-n-ten-qubit-bell.csv};
    \addplot[quimb] table[x=qubit-number, y=para-n-ten-qubit-bell-avg:quimb-simulation] {para-n-ten-qubit-bell.csv};

    \nextgroupplot[
    xlabel={Number of qubits},
    axis y line*=left,
    title={Two $N$-qubit Bell in parallel},
    mark repeat=4,
    execute at end plot visualization/.append={%
      \xdef\savedxmintnbell{\pgfkeysvalueof{/pgfplots/xmin}}%
      \xdef\savedxmaxtnbell{\pgfkeysvalueof{/pgfplots/xmax}}%
      \xdef\savedymintnbell{\pgfkeysvalueof{/pgfplots/ymin}}%
      \xdef\savedymaxtnbell{\pgfkeysvalueof{/pgfplots/ymax}}%
    },
    ]
    \addplot[fggi] table[x=qubit-number, y=para-two-n-bell-state-allein_opt:iso-fgg-total] {para-two-n-bell-state.csv};
    \addplot[fggp] table[x=qubit-number, y=para-two-n-bell-state-allein_opt:iso-perpl-fgg-total] {para-two-n-bell-state.csv};
    \addplot[qiskit] table[x=qubit-number, y=para-two-n-bell-state-avg:qiskit-simulation] {para-two-n-bell-state.csv};
    \addplot[cirq] table[x=qubit-number, y=para-two-n-bell-state-avg:qsim-simulation] {para-two-n-bell-state.csv};
    \addplot[quimb] table[x=qubit-number, y=para-two-n-bell-state-avg:quimb-simulation] {para-two-n-bell-state.csv};
  \end{groupplot}

  \begin{axis}[
    at={(b1 c1r1.south west)},
    anchor=south west,
    ymode=log,
    xmin=\savedxminhlnb,
    xmax=\savedxmaxhlnb,
    ymin=exp(\savedyminhlnb),
    ymax=exp(\savedymaxhlnb),
    hide axis,
    ]
    \addplot[qtorch] table[x=qubit-number, y=para-n-had-last-n-bell-avg:qtorch-simulation-full] {para-n-had-last-n-bell.csv};
  \end{axis}

  \begin{axis}[
    at={(b1 c2r1.south west)},
    anchor=south west,
    xmin=\savedxminphadtb,
    xmax=\savedxmaxphadtb,
    ymin=exp(\savedyminphadtb),
    ymax=exp(\savedymaxphadtb),
    ymode=log,
    hide axis,
    ]
    \addplot[qtorch] table[x=qubit-number, y=para-n-had-last-ten-qubit-bell-avg:qtorch-simulation-full] {para-n-had-last-ten-qubit-bell.csv};
  \end{axis}

  \begin{axis}[
    at={(b1 c1r2.south west)},
    anchor=south west,
    ymode=log,
    xmin=\savedxminntbell,
    xmax=\savedxmaxntbell,
    ymin=exp(\savedyminntbell),
    ymax=exp(\savedymaxntbell),
    hide axis,
    ]
    \addplot[qtorch] table[x=qubit-number, y=para-n-ten-qubit-bell-avg:qtorch-simulation-full] {para-n-ten-qubit-bell.csv};
  \end{axis}

  \begin{axis}[
    at={(b1 c2r2.south west)},
    anchor=south west,
    ymode=log,
    xmin=\savedxmintnbell,
    xmax=\savedxmaxtnbell,
    ymin=exp(\savedymintnbell),
    ymax=exp(\savedymaxtnbell),
    hide axis,
    ]
    \addplot[qtorch] table[x=qubit-number, y=para-two-n-bell-state-avg:qtorch-simulation-full] {para-two-n-bell-state.csv};
  \end{axis}


  \begin{axis}[
    at={(b1 c1r1.south west)},
    anchor=south west,
    axis x line=none,
    axis y line*=right,
    ymode=log,
    ]
    \addplot[gates] table[x=qubit-number, y=gate-number] {para-n-had-last-n-bell.csv};
  \end{axis}

  \begin{axis}[
    at={(b1 c2r1.south west)},
    anchor=south west,
    axis x line=none,
    axis y line*=right,
    ymode=log,
    ylabel={Number of gates},
    ]
    \addplot[gates] table[x=qubit-number, y=gate-number] {para-n-had-last-ten-qubit-bell.csv};
  \end{axis}

  \begin{axis}[
    at={(b1 c1r2.south west)},
    anchor=south west,
    axis x line=none,
    axis y line*=right,
    ymode=log,
    ]
    \addplot[gates] table[x=qubit-number, y=gate-number] {para-n-ten-qubit-bell.csv};
  \end{axis}

  \begin{axis}[
    at={(b1 c2r2.south west)},
    anchor=south west,
    axis x line=none,
    axis y line*=right,
    ymode=log,
    ylabel={Number of gates},
    ]
    \addplot[gates] table[x=qubit-number, y=gate-number] {para-two-n-bell-state.csv};
  \end{axis}

  \path (b1 c1r2.south west) -- (b1 c2r2.south east) coordinate[midway] (lower_middle);
  \node[anchor=north, yshift=-2.5em] at (lower_middle) {\ref{benchmark1}};
\end{tikzpicture}
\end{center}
\caption{Simulation and compilation times of our FGG-based simulators, comparing different sets of optimizations.
  For each plot, the $x$-axis shows the number of qubits, and the $y$-axis shows run time in seconds (continued).}
\label{fig:bench-2}
\end{figure}

\newpage
\pagestyle{plain}
\begin{figure}
\pgfplotsset{
  small,
  compat=1.5,
  every tick label/.append style={font=\tiny},
  label style={font=\footnotesize},
  title style={font=\footnotesize},
  ytick pos=left,
  ymode=log,
  table/col sep=comma,
  clip marker paths=false,
  colormap/BuGn-9,
  cycle list/Set1-9,
  % combine it with 'mark list*':
  cycle multiindex* list={
    mark list*\nextlist
    Set1-9\nextlist
  },
  colormap/Paired-10,
  % styles
  fggi-all/.style={
    color=RdYlGn-B,
    mark=*,
    solid,
    mark size=1.5pt,
  },
  fggi-exp/.style={
    color=RdYlGn-N,
    mark=halfcircle*,
    solid,
    mark size=1.5pt,
    mark options={rotate=90},
  },
  fggi-id/.style={
    color=Set1-B,
    mark=halfcircle*,
    solid,
    mark size=1.5pt,
    mark options={rotate=-90},
  },
  fggi-comp/.style={
    color=gray,
    mark=none,
    dotted,
    semithick,
  },
  fggp-all/.style={
    color=Set1-B,
    mark=square*,
    solid,
    mark size=1.5pt,
    mark options={rotate=45},
  },
  fggp-exp/.style={
    color=PuOr-B,
    mark=halfsquare left*,
    solid,
    mark size=2pt,
  },
  fggp-id/.style={
    color=PuOr-N,
    mark=halfsquare right*,
    solid,
    mark size=2pt,
  },
  fggp-comp/.style={
    color=YlGnBu-H,
    mark=square,
    solid,
    mark size=1.5pt,
    mark options={fill=white,rotate=45},
  },
}
\begin{center}
\begin{tikzpicture}
  \begin{groupplot}[group style={group size=2 by 3, group name=b1},
    group/horizontal sep=4em,
    group/vertical sep=5em]
    \nextgroupplot[
    xlabel={Number of qubits},
    ylabel={Running time (seconds)},
    % axis y line*=left,
    title={Hadamard to the last qubit},
    legend style={
      at={(0.5,-0.1)},
      anchor=north,legend columns=-1,
      font=\footnotesize,
      draw=none,
      legend columns=2,
      legend cell align={left},
    },
    legend to name=benchmark2,
    mark repeat=8,
    ]
    \addplot[fggi-all] table[x=qubit-number, y=had-last-qubit-allein_opt:iso-fgg-simulation] {had-last-qubit.csv};
    \addplot[fggi-exp] table[x=qubit-number, y=had-last-qubit-ein_opt:iso-fgg-simulation] {had-last-qubit.csv};
    \addplot[fggi-id] table[x=qubit-number, y=had-last-qubit-no_opt:iso-fgg-simulation] {had-last-qubit.csv};
    \addlegendentry{FGG solver with $\expand$ and $\bigl[\begin{smallmatrix}1&0\\0&1\end{smallmatrix}\bigr]$ opts};
    \addlegendentry{FGG solver with $\expand$ opt};
    \addlegendentry{FGG solver without any opt};
    \addlegendimage{fggi-comp};
    \addlegendentry{FGG compiled from ISO};

    \nextgroupplot[
    xlabel={Number of qubits},
    % axis y line*=left,
    title={Deutsch-Jozsa with the even oracle},
    mark repeat=4,
    ]
    \addplot[fggi-all] table[x=qubit-number, y=deutsch-jozsa-is-even-simplified-allein_opt:iso-fgg-simulation] {deutsch-jozsa-is-even-simplified.csv};
    \addplot[fggi-exp] table[x=qubit-number, y=deutsch-jozsa-is-even-simplified-ein_opt:iso-fgg-simulation] {deutsch-jozsa-is-even-simplified.csv};
    \addplot[fggi-id] table[x=qubit-number, y=deutsch-jozsa-is-even-simplified-no_opt:iso-fgg-simulation] {deutsch-jozsa-is-even-simplified.csv};

    \nextgroupplot[
    xlabel={Number of qubits},
    ylabel={Running time (seconds)},
    % axis y line*=left,
    title={Simon's},
    mark repeat=1,
    ]
    \addplot[fggi-all]  table[x=qubit-number, y=simon-allein_opt:iso-fgg-simulation] {simon.csv};
    \addplot[fggi-exp]  table[x=qubit-number, y=simon-ein_opt:iso-fgg-simulation] {simon.csv};
    \addplot[fggi-id]   table[x=qubit-number, y=simon-no_opt:iso-fgg-simulation] {simon.csv};

    \nextgroupplot[
    xlabel={Number of qubits},
    % axis y line*=left,
    title={Decomposed $C^{n}(X)$},
    mark repeat=2,
    ]
    \addplot[fggi-all]  table[x=qubit-number, y=mcx-allein_opt:iso-fgg-simulation] {mcx.csv};
    \addplot[fggi-exp]  table[x=qubit-number, y=mcx-ein_opt:iso-fgg-simulation] {mcx.csv};
    \addplot[fggi-id]   table[x=qubit-number, y=mcx-no_opt:iso-fgg-simulation] {mcx.csv};

    \nextgroupplot[
    xlabel={Number of qubits},
    ylabel={Running time (seconds)},
    % axis y line*=left,
    title={QFT},
    mark repeat=2,
    ]
    \addplot[fggi-all]  table[x=qubit-number, y=qft-allein_opt:iso-fgg-simulation] {qft.csv};
    \addplot[fggi-exp]  table[x=qubit-number, y=qft-ein_opt:iso-fgg-simulation] {qft.csv};
    \addplot[fggi-id]   table[x=qubit-number, y=qft-no_opt:iso-fgg-simulation] {qft.csv};

    \nextgroupplot[
    xlabel={Number of qubits},
    % axis y line*=left,
    title={Grover},
    mark repeat=2,
    ]
    \addplot[fggi-all]  table[x=qubit-number, y=grover-allein_opt:iso-fgg-simulation] {grover.csv};
    \addplot[fggi-exp]  table[x=qubit-number, y=grover-ein_opt:iso-fgg-simulation] {grover.csv};
    \addplot[fggi-id]   table[x=qubit-number, y=grover-no_opt:iso-fgg-simulation] {grover.csv};
  \end{groupplot}

  \begin{axis}[
    at={(b1 c1r1.south west)},
    anchor=south west,
    axis x line=none,
    axis y line*=right,
    ymode=log,
    ]
    \addplot[fggi-comp] table[x=qubit-number, y=had-last-qubit-avg:iso-to-fgg] {had-last-qubit.csv};
  \end{axis}

  \begin{axis}[
    at={(b1 c2r1.south west)},
    anchor=south west,
    axis x line=none,
    axis y line*=right,
    ymode=log,
    ylabel={Compile time (seconds)},
    ]
    \addplot[fggi-comp] table[x=qubit-number, y=deutsch-jozsa-is-even-simplified-avg:iso-to-fgg] {deutsch-jozsa-is-even-simplified.csv};
  \end{axis}

  \begin{axis}[
    at={(b1 c1r2.south west)},
    anchor=south west,
    axis x line=none,
    axis y line*=right,
    ymode=log,
    ]
    \addplot[fggi-comp] table[x=qubit-number, y=simon-avg:iso-to-fgg] {simon.csv};
  \end{axis}

  \begin{axis}[
    at={(b1 c2r2.south west)},
    anchor=south west,
    axis x line=none,
    axis y line*=right,
    ymode=log,
    ylabel={Compile time (seconds)},
    ]
    \addplot[fggi-comp] table[x=qubit-number, y=mcx-avg:iso-to-fgg] {mcx.csv};
  \end{axis}

  \begin{axis}[
    at={(b1 c1r3.south west)},
    anchor=south west,
    axis x line=none,
    axis y line*=right,
    ymode=log,
    ]
    \addplot[fggi-comp] table[x=qubit-number, y=qft-avg:iso-to-fgg] {qft.csv};
  \end{axis}

  \begin{axis}[
    at={(b1 c2r3.south west)},
    anchor=south west,
    axis x line=none,
    axis y line*=right,
    ymode=log,
    ylabel={Compile time (seconds)},
    ]
    \addplot[fggi-comp] table[x=qubit-number, y=grover-avg:iso-to-fgg] {grover.csv};
  \end{axis}

  \path (b1 c1r3.south west) -- (b1 c2r3.south east) coordinate[midway] (lower_middle);
  \node[anchor=north, yshift=-2.5em] at (lower_middle) {\ref{benchmark2}};
\end{tikzpicture}
\end{center}
\caption{Simulation and compilation times of our FGG-based simulators, comparing different sets of optimizations.
  For each plot, the $x$-axis shows the number of qubits, and the $y$-axis shows run time in seconds.}
\label{fig:ablation}
\end{figure}

\newpage
\pagestyle{plain}
\begin{figure}
\pgfplotsset{
  small,
  compat=1.5,
  every tick label/.append style={font=\tiny},
  label style={font=\footnotesize},
  title style={font=\footnotesize},
  ytick pos=left,
  ymode=log,
  table/col sep=comma,
  clip marker paths=false,
  colormap/BuGn-9,
  cycle list/Set1-9,
  % combine it with 'mark list*':
  cycle multiindex* list={
    mark list*\nextlist
    Set1-9\nextlist
  },
  colormap/Paired-10,
  % styles
  fggi-all/.style={
    color=RdYlGn-B,
    mark=*,
    solid,
    mark size=1.5pt,
  },
  fggi-exp/.style={
    color=RdYlGn-N,
    mark=halfcircle*,
    solid,
    mark size=1.5pt,
    mark options={rotate=90},
  },
  fggi-id/.style={
    color=Set1-B,
    mark=halfcircle*,
    solid,
    mark size=1.5pt,
    mark options={rotate=-90},
  },
  fggi-comp/.style={
    color=gray,
    mark=none,
    dotted,
    semithick,
  },
  fggp-all/.style={
    color=Set1-B,
    mark=square*,
    solid,
    mark size=1.5pt,
    mark options={rotate=45},
  },
  fggp-exp/.style={
    color=PuOr-B,
    mark=halfsquare left*,
    solid,
    mark size=2pt,
  },
  fggp-id/.style={
    color=PuOr-N,
    mark=halfsquare right*,
    solid,
    mark size=2pt,
  },
  fggp-comp/.style={
    color=YlGnBu-H,
    mark=square,
    solid,
    mark size=1.5pt,
    mark options={fill=white,rotate=45},
  },
}
\begin{center}
\begin{tikzpicture}
  \begin{groupplot}[group style={group size=2 by 3, group name=b1},
    group/horizontal sep=4em,
    group/vertical sep=5em]
    \nextgroupplot[
    xlabel={Number of qubits},
    ylabel={Running time (seconds)},
    % axis y line*=left,
    title={Bell state},
    legend style={
      at={(0.5,-0.1)},
      anchor=north,legend columns=-1,
      font=\footnotesize,
      draw=none,
      legend columns=2,
      legend cell align={left},
    },
    legend to name=benchmark3,
    mark repeat=4,
    ]
    \addplot[fggi-all] table[x=qubit-number, y=bell-state-allein_opt:iso-fgg-simulation] {bell-state.csv};
    \addplot[fggi-exp] table[x=qubit-number, y=bell-state-ein_opt:iso-fgg-simulation] {bell-state.csv};
    \addplot[fggi-id] table[x=qubit-number, y=bell-state-no_opt:iso-fgg-simulation] {bell-state.csv};
    \addlegendentry{FGG solver with $\expand$ and $\bigl[\begin{smallmatrix}1&0\\0&1\end{smallmatrix}\bigr]$ opts};
    \addlegendentry{FGG solver with $\expand$ opt};
    \addlegendentry{FGG solver without any opt};
    \addlegendimage{fggi-comp};
    \addlegendentry{FGG compiled from ISO};

    \nextgroupplot[
    xlabel={Number of qubits},
    % axis y line*=left,
    title={Two $N$-qubit Bell in cascade},
    mark repeat=4,
    ]
    \addplot[fggi-all] table[x=qubit-number, y=casc-two-n-bell-state-allein_opt:iso-fgg-simulation] {casc-two-n-bell-state.csv};
    \addplot[fggi-exp] table[x=qubit-number, y=casc-two-n-bell-state-ein_opt:iso-fgg-simulation] {casc-two-n-bell-state.csv};
    \addplot[fggi-id] table[x=qubit-number, y=casc-two-n-bell-state-no_opt:iso-fgg-simulation] {casc-two-n-bell-state.csv};

    \nextgroupplot[
    xlabel={Number of qubits},
    ylabel={Running time (seconds)},
    % axis y line*=left,
    title={Hadamard to the first qubit before Bell},
    mark repeat=1,
    ]
    \addplot[fggi-all]  table[x=qubit-number, y=casc-had-first-bell-allein_opt:iso-fgg-simulation] {casc-had-first-bell.csv};
    \addplot[fggi-exp]  table[x=qubit-number, y=casc-had-first-bell-ein_opt:iso-fgg-simulation] {casc-had-first-bell.csv};
    \addplot[fggi-id]   table[x=qubit-number, y=casc-had-first-bell-no_opt:iso-fgg-simulation] {casc-had-first-bell.csv};

    \nextgroupplot[
    xlabel={Number of qubits},
    % axis y line*=left,
    title={Hadamard to the last qubit before Bell},
    mark repeat=2,
    ]
    \addplot[fggi-all]  table[x=qubit-number, y=casc-had-last-bell-allein_opt:iso-fgg-simulation] {casc-had-last-bell.csv};
    \addplot[fggi-exp]  table[x=qubit-number, y=casc-had-last-bell-ein_opt:iso-fgg-simulation] {casc-had-last-bell.csv};
    \addplot[fggi-id]   table[x=qubit-number, y=casc-had-last-bell-no_opt:iso-fgg-simulation] {casc-had-last-bell.csv};

    \nextgroupplot[
    xlabel={Number of qubits},
    ylabel={Running time (seconds)},
    % axis y line*=left,
    title={One Hadamard in parallel with Deutsch-Jozsa},
    mark repeat=2,
    ]
    \addplot[fggi-all]  table[x=qubit-number, y=had-last-dj-even-allein_opt:iso-fgg-simulation] {had-last-dj-even.csv};
    \addplot[fggi-exp]  table[x=qubit-number, y=had-last-dj-even-ein_opt:iso-fgg-simulation] {had-last-dj-even.csv};
    \addplot[fggi-id]   table[x=qubit-number, y=had-last-dj-even-no_opt:iso-fgg-simulation] {had-last-dj-even.csv};
  \end{groupplot}

  \begin{axis}[
    at={(b1 c1r1.south west)},
    anchor=south west,
    axis x line=none,
    axis y line*=right,
    ymode=log,
    ]
    \addplot[fggi-comp] table[x=qubit-number, y=bell-state-avg:iso-to-fgg] {bell-state.csv};
  \end{axis}

  \begin{axis}[
    at={(b1 c2r1.south west)},
    anchor=south west,
    axis x line=none,
    axis y line*=right,
    ymode=log,
    ylabel={Compile time (seconds)},
    ]
    \addplot[fggi-comp] table[x=qubit-number, y=casc-two-n-bell-state-avg:iso-to-fgg] {casc-two-n-bell-state.csv};
  \end{axis}

  \begin{axis}[
    at={(b1 c1r2.south west)},
    anchor=south west,
    axis x line=none,
    axis y line*=right,
    ymode=log,
    ]
    \addplot[fggi-comp] table[x=qubit-number, y=casc-had-first-bell-avg:iso-to-fgg] {casc-had-first-bell.csv};
  \end{axis}

  \begin{axis}[
    at={(b1 c2r2.south west)},
    anchor=south west,
    axis x line=none,
    axis y line*=right,
    ymode=log,
    ylabel={Compile time (seconds)},
    ]
    \addplot[fggi-comp] table[x=qubit-number, y=casc-had-last-bell-avg:iso-to-fgg] {casc-had-last-bell.csv};
  \end{axis}

  \begin{axis}[
    at={(b1 c1r3.south west)},
    anchor=south west,
    axis x line=none,
    axis y line*=right,
    ylabel={Compile time (seconds)},
    ymode=log,
    ]
    \addplot[fggi-comp] table[x=qubit-number, y=had-last-dj-even-avg:iso-to-fgg] {had-last-dj-even.csv};
  \end{axis}

  \path (b1 c1r3.south west) -- (b1 c2r3.south east) coordinate[midway] (lower_middle);
  \node[anchor=north, yshift=-2.5em] at (lower_middle) {\ref{benchmark3}};
\end{tikzpicture}
\end{center}
\caption{Simulation and compilation times of our FGG-based simulators, comparing different sets of optimizations.
  For each plot, the $x$-axis shows the number of qubits, and the $y$-axis shows run time in seconds (continued).}
\label{fig:ablation}
\end{figure}

\newpage
\pagestyle{plain}
\begin{figure}
\pgfplotsset{
  small,
  compat=1.5,
  every tick label/.append style={font=\tiny},
  label style={font=\footnotesize},
  title style={font=\footnotesize},
  ytick pos=left,
  ymode=log,
  table/col sep=comma,
  clip marker paths=false,
  colormap/BuGn-9,
  cycle list/Set1-9,
  % combine it with 'mark list*':
  cycle multiindex* list={
    mark list*\nextlist
    Set1-9\nextlist
  },
  colormap/Paired-10,
  % styles
  fggi-all/.style={
    color=RdYlGn-B,
    mark=*,
    solid,
    mark size=1.5pt,
  },
  fggi-exp/.style={
    color=RdYlGn-N,
    mark=halfcircle*,
    solid,
    mark size=1.5pt,
    mark options={rotate=90},
  },
  fggi-id/.style={
    color=Set1-B,
    mark=halfcircle*,
    solid,
    mark size=1.5pt,
    mark options={rotate=-90},
  },
  fggi-comp/.style={
    color=gray,
    mark=none,
    dotted,
    semithick,
  },
  fggp-all/.style={
    color=Set1-B,
    mark=square*,
    solid,
    mark size=1.5pt,
    mark options={rotate=45},
  },
  fggp-exp/.style={
    color=PuOr-B,
    mark=halfsquare left*,
    solid,
    mark size=2pt,
  },
  fggp-id/.style={
    color=PuOr-N,
    mark=halfsquare right*,
    solid,
    mark size=2pt,
  },
  fggp-comp/.style={
    color=YlGnBu-H,
    mark=square,
    solid,
    mark size=1.5pt,
    mark options={fill=white,rotate=45},
  },
}
\begin{center}
\begin{tikzpicture}
  \begin{groupplot}[group style={group size=2 by 2, group name=b1},
    group/horizontal sep=4em,
    group/vertical sep=5em]
    \nextgroupplot[
    xlabel={Number of qubits},
    ylabel={Running time (seconds)},
    % axis y line*=left,
    title={Hadamard in parallel with $N$-qubit Bell},
    legend style={
      at={(0.5,-0.1)},
      anchor=north,legend columns=-1,
      font=\footnotesize,
      draw=none,
      legend columns=2,
      legend cell align={left},
    },
    legend to name=benchmark4,
    mark repeat=4,
    ]
    \addplot[fggi-all] table[x=qubit-number, y=para-n-had-last-n-bell-allein_opt:iso-fgg-simulation] {para-n-had-last-n-bell.csv};
    \addplot[fggi-exp] table[x=qubit-number, y=para-n-had-last-n-bell-ein_opt:iso-fgg-simulation] {para-n-had-last-n-bell.csv};
    \addplot[fggi-id] table[x=qubit-number, y=para-n-had-last-n-bell-no_opt:iso-fgg-simulation] {para-n-had-last-n-bell.csv};
    \addlegendentry{FGG solver with $\expand$ and $\bigl[\begin{smallmatrix}1&0\\0&1\end{smallmatrix}\bigr]$ opts};
    \addlegendentry{FGG solver with $\expand$ opt};
    \addlegendentry{FGG solver without any opt};
    \addlegendimage{fggi-comp};
    \addlegendentry{FGG compiled from ISO};

    \nextgroupplot[
    xlabel={Number of qubits},
    % axis y line*=left,
    title={Hadamard in parallel with 10-qubit Bell},
    mark repeat=4,
    ]
    \addplot[fggi-all] table[x=qubit-number, y=para-n-had-last-ten-qubit-bell-allein_opt:iso-fgg-simulation] {para-n-had-last-ten-qubit-bell.csv};
    \addplot[fggi-exp] table[x=qubit-number, y=para-n-had-last-ten-qubit-bell-ein_opt:iso-fgg-simulation] {para-n-had-last-ten-qubit-bell.csv};
    \addplot[fggi-id] table[x=qubit-number, y=para-n-had-last-ten-qubit-bell-no_opt:iso-fgg-simulation] {para-n-had-last-ten-qubit-bell.csv};

    \nextgroupplot[
    xlabel={Number of qubits},
    ylabel={Running time (seconds)},
    % axis y line*=left,
    title={$N$ 10-qubit Bell in parallel},
    mark repeat=1,
    ]
    \addplot[fggi-all]  table[x=qubit-number, y=para-n-ten-qubit-bell-allein_opt:iso-fgg-simulation] {para-n-ten-qubit-bell.csv};
    \addplot[fggi-exp]  table[x=qubit-number, y=para-n-ten-qubit-bell-ein_opt:iso-fgg-simulation] {para-n-ten-qubit-bell.csv};
    \addplot[fggi-id]   table[x=qubit-number, y=para-n-ten-qubit-bell-no_opt:iso-fgg-simulation] {para-n-ten-qubit-bell.csv};

    \nextgroupplot[
    xlabel={Number of qubits},
    % axis y line*=left,
    title={Two $N$-qubit Bell in parallel},
    mark repeat=2,
    ]
    \addplot[fggi-all]  table[x=qubit-number, y=para-two-n-bell-state-allein_opt:iso-fgg-simulation] {para-two-n-bell-state.csv};
    \addplot[fggi-exp]  table[x=qubit-number, y=para-two-n-bell-state-ein_opt:iso-fgg-simulation] {para-two-n-bell-state.csv};
    \addplot[fggi-id]   table[x=qubit-number, y=para-two-n-bell-state-no_opt:iso-fgg-simulation] {para-two-n-bell-state.csv};
  \end{groupplot}

  \begin{axis}[
    at={(b1 c1r1.south west)},
    anchor=south west,
    axis x line=none,
    axis y line*=right,
    ymode=log,
    ]
    \addplot[fggi-comp] table[x=qubit-number, y=para-n-had-last-n-bell-avg:iso-to-fgg] {para-n-had-last-n-bell.csv};
  \end{axis}

  \begin{axis}[
    at={(b1 c2r1.south west)},
    anchor=south west,
    axis x line=none,
    axis y line*=right,
    ymode=log,
    ylabel={Compile time (seconds)},
    ]
    \addplot[fggi-comp] table[x=qubit-number, y=para-n-had-last-ten-qubit-bell-avg:iso-to-fgg] {para-n-had-last-ten-qubit-bell.csv};
  \end{axis}

  \begin{axis}[
    at={(b1 c1r2.south west)},
    anchor=south west,
    axis x line=none,
    axis y line*=right,
    ymode=log,
    ]
    \addplot[fggi-comp] table[x=qubit-number, y=para-n-ten-qubit-bell-avg:iso-to-fgg] {para-n-ten-qubit-bell.csv};
  \end{axis}

  \begin{axis}[
    at={(b1 c2r2.south west)},
    anchor=south west,
    axis x line=none,
    axis y line*=right,
    ymode=log,
    ylabel={Compile time (seconds)},
    ]
    \addplot[fggi-comp] table[x=qubit-number, y=para-two-n-bell-state-avg:iso-to-fgg] {para-two-n-bell-state.csv};
  \end{axis}

  \path (b1 c1r2.south west) -- (b1 c2r2.south east) coordinate[midway] (lower_middle);
  \node[anchor=north, yshift=-2.5em] at (lower_middle) {\ref{benchmark4}};
\end{tikzpicture}
\end{center}
\caption{Simulation and compilation times of our FGG-based simulators, comparing different sets of optimizations.
  For each plot, the $x$-axis shows the number of qubits, and the $y$-axis shows run time in seconds (continued).}
\label{fig:ablation}
\end{figure}

\end{document}
